%! Setting Up the SVD — Unit 7, Lesson 7.0 — Warm-Up
\documentclass[10pt]{article}
\usepackage{linalg-article}
\usepackage{linalg-boxes}
\newcommand{\TallMath}[1]{$\displaystyle #1\rule[-1.4em]{0pt}{3.2em}$}

\begin{document}

\pageheader{Unit 7, Lesson 7.0}{Warm-Up}
\namedateperiod

\vspace{0.12in}

\begin{notesbox}{Getting Ready: Ideas We'll Use Today}
\small These three ideas from Units~1 and~4 set up today's picture of the SVD. Answer quickly.

\vspace{4pt}
\textbf{1. Matrix times a vector (Unit 1).} Compute $A\mathbf{v}$ for
$A=\begin{bmatrix} 2 & 1 \\ 1 & 2 \end{bmatrix}$ and $\mathbf{v}=\begin{bmatrix} 1 \\ 1 \end{bmatrix}$:
\[
  A\mathbf{v}=\begin{bmatrix} 2 & 1 \\ 1 & 2 \end{bmatrix}\begin{bmatrix} 1 \\ 1 \end{bmatrix}
  = \begin{bmatrix}\rule{0pt}{1.1em}\blank{1.0cm}\\[2pt]\blank{1.0cm}\end{bmatrix}.
\]

\vspace{4pt}
\textbf{2. Perpendicular vectors (Unit 4).} Two vectors are \emph{perpendicular} when their dot
product is $0$. Are $\begin{bmatrix} 1 \\ 1 \end{bmatrix}$ and $\begin{bmatrix} 1 \\ -1 \end{bmatrix}$
perpendicular?
\[
  \begin{bmatrix} 1 \\ 1 \end{bmatrix}\cdot\begin{bmatrix} 1 \\ -1 \end{bmatrix}
  = (1)(1)+(1)(-1) = \blank{1.6cm}\quad\Rightarrow\quad \text{perpendicular? }\blank{1.8cm}
\]

\vspace{4pt}
\textbf{3. Length and unit vectors (Unit 4).} The length of $\begin{bmatrix} 1 \\ 1 \end{bmatrix}$ is
$\sqrt{1^2+1^2}=\blank{1.4cm}$. To make it a \emph{unit} vector (length $1$), divide by its length:
\[
  \frac{1}{\blank{0.9cm}}\begin{bmatrix} 1 \\ 1 \end{bmatrix}.
\]
\end{notesbox}

\end{document}
